% Created 2026-06-17 Wed 23:53
% Intended LaTeX compiler: xelatex
\documentclass[11pt]{article}

\usepackage{amsmath}
\usepackage{fontspec}
\usepackage{graphicx}
\usepackage{longtable}
\usepackage{wrapfig}
\usepackage{rotating}
\usepackage[normalem]{ulem}
\usepackage{capt-of}
\usepackage{hyperref}
\usepackage{fontspec}
\setmainfont{DejaVu Serif}
\setmonofont{DejaVu Sans Mono}
\date{\today}
\title{Relatorio}
\hypersetup{
 pdfauthor={},
 pdftitle={Relatorio},
 pdfkeywords={},
 pdfsubject={},
 pdfcreator={},
 pdflang={English}}
\begin{document}

\maketitle
\tableofcontents

\section{Sobre a organização dos arquivos:}
\label{sec:orgc13b5e0}


Os principais arquivos são esses:
A estratégia foi agrupar atividades semelhantes em cada arquivo.
Cada header file (.h) incluí o protótipo das funções declaradas no respectivo source file (.c)

\begin{verbatim}
tree | grep -e '.c' -e '.h'
\end{verbatim}

\begin{center}
\begin{tabular}{l l p{10cm}}
├── & game.c & Definição das estruturas e funções que representam inimigos, jogador, sua movimentação, seus status\\
├── & game.h & \\
├── & main.c & Arquivo principal, onde se definem as variáveis e estruturas além do game-loop\\
├── & mapa.c & Trata do vetor de char que representa a posição inicial de todos elementos do mapa, a partir dos arquivos mapaX.txt\\
├── & mapa.h & \\
├── & menu.c & Trata dos diferentes menus e suas funcionalidades, incluindo suas opções e aparência.\\
├── & menu.h & \\
├── & placar.bin & Guarda os 10 melhores scores\\
├── & render.c & Lida com a renderização do mapa e todas entidades.\\
└── & render.h & \\
\end{tabular}
\end{center}
\section{Funcionamento}
\label{sec:org3c737ed}
\subsection{Estruturas}
\label{sec:org3ca45d1}
De um modo geral, foram definidas estruturas que representam inimigos (flames), o jogador (player) além de um tipo base (base);
Essas estruturas contém atributos basicos como posição, velocidade além de flags como ``está vivo?'', ``está pindurado?'';
Essas estruturas são utilizadas majoritariamente na implementação da jogabilidade do game-loop, isto é, são utilizadas na renderização das entidades (inimigos e jogador) no mapa durante o jogo, bem como sua movimentação, interações entre si (colisão entre entidades, por exemplo, matar o player caso toque em um inimigo) e entre o mapa (Logica que previne entidades de penetrar as plataformas).
Além dessas (que são as principais), foram utilizadas extruturas auxiliares: ``entities'' que funciona como vetor contendo todas entidades ativas, e a estrutura ``base'';
Essa ultima contem todos os campos que são comuns a jogadores e inimigos (principalmente referentes a posição e locomoção):
Todos os campos comuns a essas tres (inimigos, jogador e base) são declarados na mesma ordem e no inicio de cada tipo;
Isso permite que criemos funções mais genericas, que recebam de parametro tanto inimigos ou jogadores, utilizando um cast para esse tipo base.
\begin{itemize}
\item Por exemplo, temos a função
void calculaCantosInt2(base* entity);
\end{itemize}
podemos chama-la para atualizar os cantos (que são campos dos jogadores e inimigos) de entidades:
\begin{verbatim}
        calculaCantosInt2((base*)jogador);
            ou
        calculaCantosInt2((base*)inimigo);
\end{verbatim}
\subsection{Game-loop}
\label{sec:orga5620a3}
Dentro de main.c, o game-loop se divide em algumas subpartes, cada uma sendo um 'case' em um switch case principal:
\subsubsection{São 5 casos:}
\label{sec:org91ff660}
\begin{enumerate}
\item 'case 0' define a representação e funcionalidade do menu inicial;
\item 'case 1' define a principal parte do loop, que trata do 'jogo rodando';
É aqui que o jogador pode se locomover, saltar, ser morto pelos inimigos\ldots{}
\item 'case 2' faz funcionar a tela ``você morreu'', quando o jogador, more\ldots{}
\item 'case 3' a tela de transição quando o jogador passa de fase
\item 'case 4' a tela de Pause;
\end{enumerate}

\begin{itemize}
\item Obs: os 'case's 0 e 4 apresentam menus interativos, enquanto 2 e 3 apenas mostram mensagens;
\item A cada game-loop, uma flag chamada 'gameMode' decide qual dos switch cases será executado, sendo ela alterada apropriadamente em cada situação.
\end{itemize}
\subsection{Um pouco mais de perto (cases)}
\label{sec:org6410a70}
\subsubsection{Renderização}
\label{sec:org0505d4e}
De uma forma bem geral, em cada um dos casos, a ultima coisa a ser feita é a renderização, que ocorre no escopo de ``Drawing()'';
Invariavelmente, isso segue a lógica de primeiramente, tornar tudo preto (ClearBackground(BLACK);) e depois se desenha o que deve ser visto; Isso se repete todo frame (60 vezes por segundo)
Ex:
\begin{verbatim}
          BeginDrawing();
              ClearBackground(BLACK);
              menuDraw(1,menu,ptrStatus);
          EndDrawing();
\end{verbatim}
Isso é verdade em todos os casos; Uma função decide o que será desenhado, (no caso acima ``menuDraw(1,menu,ptrStatus);'')
\begin{itemize}
\item \textbf{Tangente, render.c}
\end{itemize}
Apesar do esforço para separar o código em arquivos de semelhantes funcionalidades, render.c acaba lidando apenas com uma função (e bem simples), responsável pela renderização no caso 'case 1' de jogo; No entanto, o restante das funções de renderização ficam localizadas no menu.c.
\subsubsection{O principal (case 1)}
\label{sec:org2acdec3}
Esse é de longe o caso mais complexo.
Ele pode ser subdividido em trechos:
\begin{enumerate}
\item 'ouvir' por inputs / mover entidades
De um modo geral, nessa parte inicial, são registrados inputs via teclado, tanto sobre movimentação quanto para pausar o jogo.
No primeiro caso, geralmente, atualiza-se o struct do jogador apropriadamente (ex. sua posição horizontal é acrescida caso o jogador aperte 'D'). Outras atualizações acometem todas entidades nesse estado, como a gravidade que altera os campos de velocidade vertical de todas entidades presentes, além da IA dos inimigos, que altera a posição dos inimigos\ldots{}
No segundo caso, caso o jogador aperte 'P', a flag gameMode é alterada, de modo que no frame seguinte o estado do jogo passa para o switch case apropriado.
\item calcular colisões e checar outros eventos
Em geral, isso se dá checando se a entidade em questão ``invadiu'' os limites de um bloco sólido (as plataformas por exemplo), teleportando tal entidade de volta; Como a logica detecta essa invasão antes de renderizar (além do fato de serem 60 FPS), ocorre a ilusão de que o mapa é sólido;
Outro caso semelhante é o das escadas, em que detecta-se a proximidade da entidade com o inicio ou fim das escadas, ou do jogador com o Fim da fase.
\item Renderizar
\end{enumerate}
\section{Sobre como usar:}
\label{sec:org47b97fe}
\subsection{Utilizando bash:}
\label{sec:org4b1c8e3}
\begin{enumerate}
\item Clonando o repositório github:
\begin{verbatim}
cd
git clone https://github.com/wil3212/InfDK.git
cd InfDk
\end{verbatim}

\item Uma vez dentro do diretório basta rodar o arquivo ``jogo'':
\end{enumerate}
\begin{verbatim}
./jogo
\end{verbatim}

\begin{enumerate}
\item Para copilar e rodar:
mas requer que raylib esteja instalada
\end{enumerate}
\begin{verbatim}
make -k && ./jogo
\end{verbatim}

\textbf{Obs:}
\begin{itemize}
\item É importante que o terminal siga aberto para que o jogo funcione, fechar o terminal causará o fim do jogo.
\item Atualmente, ainda não foi implementada a visualização do placar dentro da GUI do jogo (na opção 'Ranking'), no entanto, ao encerrar uma jogada (quando acabam-se as vidas), o jogador deve informar seu nome no terminal e apertar enter (o nome deve ter até 20 caracteres), ficando salvo a posição do jogador no ranking;
\item A opção 'Ranking' printa os valores no terminal;
\end{itemize}
\subsection{Jogabilidade:}
\label{sec:org71391a8}
\begin{itemize}
\item Navega pelos menus com 'W' e 'D', selecionando a opção com Enter ou Espaço;
\item Em jogo: Anda para esquerda e direita com 'A'e 'D', pula com Espaço, sobe e desce escadas com 'W' e 'S'
\item Pausa o jogo com TAB ou 'P', despausa precionando novamente ou selecionando a opção no menu
\end{itemize}
\end{document}
